\documentclass[12pt]{article}
\usepackage[utf8]{inputenc}
\usepackage[T1]{fontenc}
\usepackage{amsmath, amssymb, geometry, booktabs}
\usepackage{enumitem}
\geometry{margin=1in}
\setlength{\parindent}{0pt}
\renewcommand{\arraystretch}{1.2}

\begin{document}



\section*{Problem 1 (13 marks)}

\textbf{Given:} Home $(K,L)=(300,900)$ so $(K/L)_H = 300/900 = 0.33$.\\
Foreign $(K,L)=(800,800)$ so $(K/L)_F = 800/800 = 1.00$.\\
$T$ is labour-intensive; $S$ is capital-intensive. Identical technologies; homothetic preferences; perfect competition; full employment.

\begin{enumerate}[label=(\arabic*)]
\item \textbf{Relative factor abundance.} A country is \emph{labour-abundant} if it has a lower $K/L$ than its trading partner. Since $0.33 < 1.00$, Home is labour-abundant and Foreign is capital-abundant.

\item \textbf{Trade pattern (H--O theorem).} With identical technologies and homothetic preferences, each country exports the good that uses its abundant factor intensively. Therefore, Home (labour-abundant) exports $T$; Foreign (capital-abundant) exports $S$.

\item \textbf{Factor prices (Stolper--Samuelson).} Opening to trade raises the relative price of the export good in each country (versus autarky). Hence in Home, the relative price of $T$ rises $\Rightarrow$ the \emph{real wage} (return to labour, the intensive factor in $T$) \textbf{rises} and the \emph{real rental rate} \textbf{falls}. In Foreign, the relative price of $S$ rises $\Rightarrow$ the \emph{real rental} \textbf{rises} and the \emph{real wage} \textbf{falls}. (Results hold in real terms with respect to both goods under the usual assumptions.)

\item \textbf{A 10\% rise in $P_T/P_S$ (Home).} By the magnification effect, the percentage change in factor prices exceeds the percentage change in the price of the good that uses that factor intensively. Thus the \emph{wage} rises by \(>\!10\%\) (in real terms), while the \emph{rental} falls (in real terms). Sectoral outputs adjust in the long run: more resources flow to $T$ and away from $S$ (Rybczynski logic at given prices; with the price change, $T$ expands and $S$ contracts).

\item \textbf{Mechanism for factor price equalization (FPE).} With identical technologies, common goods prices under free trade pin down unit-cost equations:
\[
\begin{aligned}
P_T &= a_{LT}w + a_{KT}r,\\
P_S &= a_{LS}w + a_{KS}r.
\end{aligned}
\]
Given $(P_T,P_S)$ and strictly concave, competitive technology, $(w,r)$ is determined uniquely. If both countries produce both goods (diversification) and share technologies, they face the same $(P_T,P_S)$ and hence converge to the same $(w,r)$ in the \emph{cone of diversification} (the FPE set).
\end{enumerate}

\hrule
\vspace{0.6em}
\newpage
\section*{Problem 2 (13 marks)}

\textbf{Given:} PPF $A^2 + 2E^2 = 900$. World relative price $P_E/P_A = 3$. Preferences satisfy $E = 0.5\,A$ at the consumption point.

\medskip
\textbf{Normalization.} Set $P_A = 1 \Rightarrow P_E = 3$. National income at production point: $V = A + 3E$.

\subsection*{(1) Value-maximizing production $(A_P,E_P)$}
Maximize $V = A + 3E$ subject to $A^2 + 2E^2 = 900$. At an interior optimum, the slope of the PPF equals the slope of the isovalue line:

\[
\text{PPF: } A^2+2E^2=900 \Rightarrow 2A\,dA + 4E\,dE = 0 \Rightarrow \frac{dE}{dA} = -\frac{A}{2E}.
\]
\[
\text{Isovalue slope} = -\frac{P_A}{P_E} = -\frac{1}{3}.
\]
Tangency gives $-\dfrac{A}{2E} = -\dfrac{1}{3} \Rightarrow \dfrac{A}{2E} = \dfrac{1}{3} \Rightarrow 3A = 2E \Rightarrow E = \frac{3}{2}A.$

Substitute into PPF:
\[
A^2 + 2\!\left(\frac{3}{2}A\right)^{\!2} = 900
\;\Rightarrow\;
A^2 + 2\cdot \frac{9}{4}A^2 = \left(1+\frac{9}{2}\right)A^2 = \frac{11}{2}A^2 = 900.
\]
Hence
\[
A_P^2 = \frac{1800}{11},\quad
A_P = \sqrt{\frac{1800}{11}},\quad
E_P = \frac{3}{2}A_P.
\]
It is often convenient to keep the exact form. Numerically,
\[
A_P \approx 12.788,\qquad E_P \approx 19.182.
\]

\subsection*{(2) Consumption $(A_C,E_C)$, export good, and trade values}
Production value (national income) at world prices:
\[
V = A_P + 3E_P = A_P + 3\!\left(\frac{3}{2}A_P\right) = \frac{11}{2}A_P.
\]
Budget line at world prices: $A + 3E = V$. Preferences require $E=0.5\,A$. Substitute into budget:
\[
A + 3(0.5A) = 2.5A = V \;\Rightarrow\; A_C = \frac{V}{2.5} = \frac{\frac{11}{2}A_P}{\frac{5}{2}} = \frac{11}{5}A_P,\quad
E_C = 0.5\,A_C = \frac{11}{10}A_P.
\]
Thus
\[
A_C = \frac{11}{5}A_P,\qquad E_C = \frac{11}{10}A_P.
\]
Trade pattern (compare $P$ vs.\ $C$):
\[
\text{Exports of }E:\; X_E = E_P - E_C = \left(\frac{3}{2} - \frac{11}{10}\right)A_P = \frac{4}{10}A_P = \frac{2}{5}A_P > 0,
\]
\[
\text{Imports of }A:\; M_A = A_C - A_P = \left(\frac{11}{5}-1\right)A_P = \frac{6}{5}A_P > 0.
\]
So the country \textbf{exports electronics} and \textbf{imports agricultural products}.

Values at world prices:
\[
\text{Export value} = P_E X_E = 3\cdot \frac{2}{5}A_P = \frac{6}{5}A_P,\qquad
\text{Import value} = P_A M_A = 1\cdot \frac{6}{5}A_P = \frac{6}{5}A_P.
\]

\subsection*{(3) Balanced trade verification}
Export and import values are equal at world prices ($\frac{6}{5}A_P$), so trade is balanced:
\[
P_E X_E = P_A M_A.
\]

\subsection*{(4) Effect of an increase in $P_E/P_A$ on production, consumption, and welfare}
An increase in $P_E/P_A$ (steeper world price line) implies:
\begin{itemize}
\item \textbf{Production:} Tangency condition $-\frac{A}{2E} = -\frac{P_A}{P_E}$ gives a more negative slope on the right-hand side $\Rightarrow$ production tilts toward $E$ (higher $E/A$).
\item \textbf{Consumption:} The trade (budget) line through the new production point pivots outward in $E$--intercepts; with homothetic preferences, consumption shifts toward more $E$ as well.
\item \textbf{Welfare:} Because the country is an $E$-exporter, a rise in $P_E/P_A$ improves its \emph{terms of trade}. With no distortions, the budget set expands and the country reaches a higher indifference curve. \textbf{Welfare rises.}
\end{itemize}

\subsection*{(5) Export-biased growth concentrated in electronics}
\begin{itemize}
\item \textbf{Relative supply:} The PPF rotates outward more in the $E$ direction; at given prices, $RS=Q_E/Q_A$ rises (rightward shift).
\item \textbf{Large-country effect on prices:} If the country is large, the world relative price $P_E/P_A$ tends to \emph{fall} (worsening the country’s terms of trade) because its export good becomes more abundant globally.
\item \textbf{Welfare:} Two opposing effects --- (i) \emph{production gain} from a larger PPF; (ii) \emph{terms-of-trade loss}. Typically welfare still rises; only in extreme cases with very inelastic foreign demand could the TOT deterioration outweigh the production gain (``immiserizing growth'').
\end{itemize}

\bigskip
\textbf{Numerical illustration (optional).} Using $A_P\approx 12.788$:
\[
E_P \approx 19.182,\quad V=\frac{11}{2}A_P\approx 70.334,\quad
A_C=\frac{11}{5}A_P\approx 28.134,\quad E_C=\frac{11}{10}A_P\approx 14.067.
\]
Exports: $X_E\approx 5.115$ with value $3X_E\approx 15.345$; Imports: $M_A\approx 15.346$ with value $\approx 15.346$ (balanced up to rounding).

\end{document}
